% Блок-схема программной архитектуры (русская версия рис. 3.2)
\begin{tikzpicture}[
    font=\small,
    box/.style={draw, rounded corners=2pt, align=center, minimum height=8mm, inner sep=3pt, fill=white},
    port/.style={draw, thick, align=center, minimum width=16mm, minimum height=7mm, fill=gray!15},
    ns/.style={draw, dashed, rounded corners=4pt, inner sep=5pt},
    lbl/.style={font=\footnotesize\itshape, align=center},
    >=Stealth]

    \node[box, minimum width=38mm] (gui) at (0,0) {управляющее приложение\\(Tkinter, дисплей)};
    \node[box, minimum width=30mm] (sluzby) at (5.2,0.6) {службы systemd:\\порты, состояние\\портов, питание};
    \node[box, minimum width=30mm] (vysledky) at (5.2,-0.8) {каталог результатов\\JSON, CSV, pcap};
    \node[box, minimum width=18mm] (wlan) at (-4.4,0) {\texttt{wlan0}\\Wi-Fi};
    \node[ns, fit=(gui)(sluzby)(vysledky)(wlan)] (vychozi) {};
    \node[lbl, anchor=south west] at (vychozi.north west) {основное пространство имён};

    \node[box, minimum width=40mm] (skripty1) at (-3.2,-3.6) {измерительные скрипты\\(анализатор, ping, iperf3, аудит)};
    \node[port] (mon0) at (-3.2,-5.1) {\texttt{mon0}};
    \node[ns, fit=(skripty1)(mon0)] (nsmon) {};
    \node[lbl, rotate=90, anchor=south] at (nsmon.west) {\texttt{analyzer\_monitor}};

    \node[box, minimum width=40mm] (skripty2) at (3.2,-3.6) {измерительные скрипты\\(генератор, ping, iperf3, аудит)};
    \node[port] (snd0) at (3.2,-5.1) {\texttt{snd0}};
    \node[ns, fit=(skripty2)(snd0)] (nssnd) {};
    \node[lbl, rotate=-90, anchor=south] at (nssnd.east) {\texttt{analyzer\_sender}};

    \coordinate (rozcesti) at (0,-1.9);
    \draw[<->] (gui.south) -- (rozcesti) -| (skripty1.north);
    \draw[<->] (rozcesti) -| (skripty2.north);
    \node[lbl, fill=white] at (rozcesti) {запуск командой \texttt{ip netns exec}\\и возврат строки \texttt{RESULT}};
    \draw[->] (sluzby.west) -- node[lbl, above] {состояние} (sluzby.west -| gui.east);
    \draw[->] (vysledky.west -| gui.east) -- node[lbl, below] {запись} (vysledky.west);

    \draw[thick] (mon0.south) -- ++(0,-0.6) -| node[lbl, pos=0.25, below] {тестируемый сегмент или перемычка между портами} (snd0.south);
\end{tikzpicture}
